\subsection{Multiplication endomorphisms}

Let X be a locally compact topological space and let \(\mu\) be a positive measure on X. Let \(p \in [1, +\infty[\) .

Let \(g \in \mathcal{L}^{\infty}(\mathrm{X},\mu)\) . Denote by \(m_{g}\) the map \(f \mapsto g \cdot f\) of \(\mathcal{L}^{p}(\mathrm{X},\mu)\) into itself. The map \(m_{g}\) is linear and continuous since

\[
\mathrm{N} _ {p} (m _ {g} (f)) \leqslant \mathrm{N} _ {\infty} (g) \mathrm{N} _ {p} (f)\tag{1}
\]

for every function \(f \in \mathcal{L}^{p}(\mathrm{X}, \mu)\) . In particular, the function \(m_{g}(f)\) is \(\mu\) -negligible if f is \(\mu\) -negligible. The map \(m_{g}\) therefore induces, by passing to quotients, an endomorphism of \(\mathrm{L}^{p}(\mathrm{X}, \mu)\) , which will be denoted by \(\widetilde{m}_{g}\) .

Lemma 7. --- Let \(g \in \mathcal{L}^{\infty}(\mathrm{X},\mu)\) . The endomorphism \(\tilde{m}_g\) of \(\mathrm{L}^{p}(\mathrm{X},\mu)\) is injective if and only if the set of \(x \in \mathrm{X}\) such that \(g(x) = 0\) is locally \(\mu\) -negligible.

Let A denote the set of \(x \in X\) such that \(g(x) = 0\) .

Suppose that A is locally \(\mu\) -negligible. Let \(f \in \mathcal{L}^p(\mathrm{X}, \mu)\) and \(\widetilde{f}\) its class in \(\mathrm{L}^p(\mathrm{X}, \mu)\) . Suppose that \(\widetilde{m}_g(\widetilde{f}) = 0\) . By definition, this is the case if and only if the function \(fg\) is \(\mu\) -negligible, so that the set B of \(x \in \mathrm{X}\) such that \(f(x)g(x) \neq 0\) is \(\mu\) -negligible. The function \(f\) vanishes outside \(\mathrm{A} \cup \mathrm{B}\) , hence is \(\mu\) -negligible. This implies that the endomorphism \(\widetilde{m}_g\) is injective.

Conversely, suppose that A is not locally \(\mu\) -negligible. Let C be a compact subset of X such that A∩C is not \(\mu\) -negligible, and let \(\varphi\) be the characteristic function of A∩C. The class of the function \(\varphi\) in L \(^{p}\) (X, \(\mu\) ) is not zero, but that of \(m_{g}(\varphi)\) is, so \(m_{g}\) is not injective.

Since the product of a locally \(\mu\) -negligible function and a \(\mu\) -negligible function is \(\mu\) -negligible, the endomorphism \(\widetilde{m}_g\) depends only on the class of \(g\) in \(\mathrm{L}^{\infty}(\mathrm{X},\mu)\) . For \(\widetilde{g}\in \mathrm{L}^{\infty}(\mathrm{X},\mu)\) , one also denotes by \(\widetilde{m}_{\widetilde{g}}\) the endomorphism \(\widetilde{m}_g\) of \(\mathrm{L}^p (\mathrm{X},\mu)\) for every function \(g\in \mathcal{L}^{\infty}(\mathrm{X},\mu)\) of which \(\widetilde{g}\) is the class. One says that \(\widetilde{m}_{\widetilde{g}}\) is the multiplication endomorphism by \(\widetilde{g}\) in \(\mathrm{L}^p (\mathrm{X},\mu)\) .

PROPOSITION 5. --- Let \(p \in [1, +\infty[\) . The map \(m: g \mapsto \widetilde{m}_g\) from \(\mathrm{L}^{\infty}(\mathrm{X}, \mu)\) into \(\mathcal{L}(\mathrm{L}^{p}(\mathrm{X}, \mu))\) is an isometric morphism of unital Banach algebras.

For \(a\) , \(b\) in \(\mathbf{C}\) and \(g_{1}\) , \(g_{2}\) in \(\mathrm{L}^{\infty}(\mathrm{X},\mu)\) , one immediately verifies that \(\widetilde{m}_{ag_1+bg_2}=a\widetilde{m}_{g_1}+b\widetilde{m}_{g_2}\) , so the map \(\widetilde{m}\) is linear. Moreover, one has \(\widetilde{m}_{1}=1_{\mathrm{L}^{p}(\mathrm{X},\mu)}\) and \(\widetilde{m}_{g_{1}g_{2}}=\widetilde{m}_{g_{1}}\widetilde{m}_{g_{2}}\) , so the map \(\widetilde{m}\) is a unital algebra morphism from \(\mathrm{L}^{\infty}(\mathrm{X},\mu)\) into \(\mathcal{L}(\mathrm{L}^{p}(\mathrm{X},\mu))\) .

Formula (1) above shows that \(\widetilde{m}\) has norm \(\leqslant 1\) .

Let \(g \in \mathcal{L}^{\infty}(\mathrm{X},\mu)\) and \(\widetilde{g}\) be its class in \(\mathrm{L}^{\infty}(\mathrm{X},\mu)\) . Let \(\varepsilon > 0\) . The set Y of \(x \in X\) such that \(|g(x)| > \| \widetilde{g}\|_{\infty} - \varepsilon\) is not locally \(\mu\) -negligible. There therefore exists a compact subset C of X such that \(\mu (\mathrm{Y} \cap \mathrm{C}) > 0\) . Let \(\varphi\) be the characteristic function of \(\mathrm{Y} \cap \mathrm{C}\) . Since

\[
\begin{array}{c} \| \widetilde {m} _ {\widetilde {g}} (\varphi) \| _ {p} = \Bigl (\int_ {\mathrm{X}} | \varphi g | ^ {p}   d \mu \Bigr) ^ {1 / p} \geqslant \Bigl (\int_ {\mathrm{Y}} | \varphi g | ^ {p}   d \mu \Bigr) ^ {1 / p} \\ \geqslant (\| \widetilde {g} \| _ {\infty} - \varepsilon) \mu (\mathrm{Y} \cap \mathrm{C}) ^ {1 / p} = (\| \widetilde {g} \| _ {\infty} - \varepsilon) \| \varphi \| _ {p}, \end{array}
\]

it follows that \(\|\widetilde{m}_{\widetilde{g}}\|\geqslant\|\widetilde{g}\|_{\infty}-\varepsilon\) . Since \(\varepsilon\) is arbitrary, we deduce that \(\|\widetilde{m}_{\widetilde{g}}\|\geqslant\|\widetilde{g}\|_{\infty}\) , which proves that \(\widetilde{m}\) is isometric.

Lemma 8. --- Let \((g_n)\) be a bounded sequence in \(\mathcal{L}^\infty(\mathrm{X},\mu)\) converging simply \(\mu\) -almost everywhere. Let \(\widetilde{m} \in \mathrm{L}^\infty(\mathrm{X},\mu)\) be the class of its limit. Then \(\widetilde{m}_{g_n}\) converges to \(\widetilde{m}_{\widetilde{g}}\) in the space \(\mathcal{L}(\mathrm{L}^p(\mathrm{X},\mu))\) endowed with the topology of simple convergence.

Let \(f \in \mathcal{L}^p(\mathrm{X},\mu)\) . The sequence \((g_{n}f)\) is bounded in \(\mathcal{L}^p(\mathrm{X},\mu)\) and converges simply \(\mu\) -almost everywhere to \(gf\) . By the Lebesgue theorem (INT, IV, p. 137, § 3, no. 7, th. 6), the sequence \((g_{n}f)\) converges to \(gf\) in \(\mathcal{L}^p(\mathrm{X},\mu)\) , and the assertion follows.

We now consider the case p = 2.

PROPOSITION 6. --- The map m: \(g \mapsto \widetilde{m}_{g}\) is a unital isometric morphism of the stellar algebra \(\mathrm{L}^{\infty}(\mathrm{X}, \mu)\) into the stellar algebra \(\mathcal{L}(\mathrm{L}^{2}(\mathrm{X}, \mu))\) .

In particular, for every \(g \in \mathrm{L}^{\infty}(\mathrm{X}, \mu)\) the multiplication endomorphism \(\widetilde{m}_{g}\) is normal; it is Hermitian if and only if g is locally real-valued \(\mu\) -almost everywhere.

By prop. 5, the map \(\tilde{m}\) is an injective, isometric morphism of unital Banach algebras from \(\mathrm{L}^{\infty}(\mathrm{X},\mu)\) into \(\mathcal{L}(\mathrm{L}^{2}(\mathrm{X},\mu))\) . Let \(g\in \mathcal{L}^{\infty}(\mathrm{X},\mu)\) . For \(f_{1}\) and \(f_{2}\in \mathcal{L}^{2}(\mathrm{X},\mu)\) , we have

\[
\langle f _ {1} \mid \widetilde {m} _ {g} (f _ {2}) \rangle = \int_ {\mathrm{X}} \overline {{f _ {1} (x)}} g (x) f _ {2} (x) d \mu (x) = \langle \widetilde {m} _ {\overline {{g}}} (f _ {1}) \mid f _ {2} \rangle ,
\]

from which it follows that \(\widetilde{m}_g^* = \widetilde{m}_{\overline{g}}\) , which proves that \(m\) is an involutive morphism. The last assertions follow from this (cf. I, p. 106, prop. 5).

COROLLARY. --- Let \(g \in \mathrm{L}^{\infty}(\mathrm{X},\mu)\) .

\begin{enumerate}
\def\labelenumi{\alph{enumi})}
\item
  The spectrum of \(\widetilde{m}_g\) in \(\mathcal{L}(\mathrm{L}^2 (\mathrm{X},\mu))\) is the \(\mu\) -essential image of \(g\) ;
\item
  The endomorphism \(\widetilde{m}_g\) is positive if and only if \(g\) is locally positive \(\mu\) -almost everywhere;
\item
  For any function \(f \in \mathcal{C}(\mathrm{Sp}(\widetilde{m}_g))\) , we have \(f(\widetilde{m}_g) = \widetilde{m}_{f \circ g}\) .
\end{enumerate}

Since \(\widetilde{m}\) is injective, we have \(\mathrm{Sp}(\widetilde{m}_{g}) = \mathrm{Sp}(g)\) by prop. 5 of I, p. 106; the result then follows from lemma 6 of IV, p. 185, proposition 4 of IV, p. 185, and definition 6 of I, p. 115.

PROPOSITION 7. --- The image of the morphism \(\tilde{m}\) from \(\mathrm{L}^{\infty}(\mathrm{X},\mu)\) into \(\mathcal{L}(\mathrm{L}^{2}(\mathrm{X},\mu))\) is a maximal commutative subalgebra of the algebra \(\mathcal{L}(\mathrm{L}^{2}(\mathrm{X},\mu))\) .

It suffices to prove that if \(u \in \mathcal{L}(\mathrm{L}^{2}(\mathrm{X}, \mu))\) is an endomorphism that commutes with \(\widetilde{m}_{g}\) for every function \(g \in \mathcal{L}^{\infty}(\mathrm{X}, \mu)\) , then u belongs to the image of \(\widetilde{m}\) . Let u be such an endomorphism.

We denote \(\tilde{f}\) the class in \(\mathrm{L}^2 (\mathrm{X},\mu)\) of a function \(f\in \mathcal{L}^2 (\mathrm{X},\mu)\) . Denote by \(v\) the linear map from \(\mathcal{L}^2 (\mathrm{X},\mu)\) into \(\mathrm{L}^2 (\mathrm{X},\mu)\) defined by \(f\mapsto u(\widetilde{f})\) . We have \(v\circ m_g = \widetilde{m}_g\circ v\) for all \(g\in \mathcal{L}^{\infty}(\mathrm{X},\mu)\) .

For any subset \(\mu\) -measurable Y of X, we denote \(\varphi_{\mathrm{Y}}\) its characteristic function; the endomorphism \(\widetilde{m}_{\varphi_{\mathrm{Y}}}\) is an orthogonal projector of \(\mathrm{L}^2 (\mathrm{X},\mu)\) .

Suppose first that X is compact, so that the class of the constant function 1 belongs to \(\mathrm{L}^{2}(\mathrm{X},\mu)\) . Let g be a function in \(\mathcal{L}^{2}(\mathrm{X},\mu)\) whose class in \(\mathrm{L}^{2}(\mathrm{X},\mu)\) is \(v(1)\) .

Let \(c\) be a positive real number and Y the set of \(x \in \mathrm{X}\) such that \(|g(x)| \geqslant c\) ; it is a \(\mu\) -measurable set. We have in \(\mathrm{L}^2(\mathrm{X}, \mu)\) the equalities \(\widetilde{\varphi_{\mathrm{Y}}} g = \widetilde{m}_{\varphi_{\mathrm{Y}}}(v(1)) = v(m_{\varphi_{\mathrm{Y}}}(1)) = v(\varphi_{\mathrm{Y}})\) .

Since moreover \(|c\varphi_{\mathrm{Y}}| \leqslant |g\varphi_{\mathrm{Y}}|\) , we obtain

\[
c ^ {2} \mu (\mathrm{Y}) = \int_ {\mathrm{X}} (c \varphi_ {\mathrm{Y}}) ^ {2} d \mu \leqslant \int_ {\mathrm{X}} | \varphi_ {\mathrm{Y}} g | ^ {2} d \mu = \| \widetilde {u} (\varphi_ {\mathrm{Y}}) \| ^ {2} \leqslant \| u \| ^ {2} \mu (\mathrm{Y}).
\]

This inequality implies that \(\mu(\mathrm{Y}) = 0\) if \(c > \|u\|\) , so that the function \(g\) is bounded \(\mu\) -almost everywhere by \(\|u\|\) . For every function \(f \in \mathcal{C}(\mathrm{X})\) , we finally have

\[
u (\widetilde {f}) = v (m _ {f} (1)) = \widetilde {m} _ {f} (v (1)) = \widetilde {m} _ {f} (\widetilde {g}) = \widetilde {m} _ {g} (\widetilde {f}),
\]

whence the equality \(u = \tilde{m}_{g}\) , and in particular \(\mathrm{N}_{\infty}(g) = \| u\|\) .

Consider now the general case. There exists a locally countable family \((\mathrm{X}_{i})_{i\in\mathrm{I}}\) of pairwise disjoint compact subsets of X such that \(Z = X - \bigcup_{i \in I} X_i\) is locally \(\mu\) -negligible (INT, IV, p.190, § 5, no. 9, prop. 14). Let \(\mu_i\) be the measure induced by \(\mu\) on \(X_i\) (INT, IV, p. 186, § 5, no. 7, def. 4).

By prop. 1 of IV, p. 179, for every \(i \in I\) the image \(\mathrm{E}_i\) of \(\widetilde{m}_{\varphi_{\mathrm{X}_i}}\) is identified with \(\mathrm{L}^2(\mathrm{X}_i, \mu_i)\) by \(f \mapsto f|\mathrm{X}_i\) . For every function \(g \in \mathcal{L}^\infty(\mathrm{X}, \mu)\) the multiplication endomorphism \(\widetilde{m}_g\) commutes with \(\widetilde{m}_{\varphi_{\mathrm{X}_i}}\) , hence leaves stable \(\mathrm{E}_i\) , and the endomorphism of \(\mathrm{E}_i\) deduced from \(\widetilde{m}_g\) by passage to the subspaces coincides with the endomorphism of multiplication by \(g|\mathrm{X}_i\) on \(\mathrm{L}^2(\mathrm{X}_i, \mu_i)\) .

Since \(u\) commutes with \(m_{\varphi_{X_i}}\) , it also leaves stable the subspace \(\mathrm{E}_i\) . Denote by \(u_i\) the endomorphism of \(\mathrm{L}^2 (\mathrm{X}_i,\mu_i)\) deduced from \(u\) by passage to the subspaces.

The map from \(\mathcal{L}^{\infty}(\mathrm{X},\mu)\) into \(\mathcal{L}^{\infty}(\mathrm{X}_{i},\mu_{i})\) defined by \(g\mapsto g|\mathrm{X}_{i}\) being surjective, the hypothesis implies that \(u_{i}\) commutes with \(\widetilde{m}_{g}\) for every function \(g\in \mathcal{L}^{\infty}(\mathrm{X}_{i},\mu_{i})\) . By the first part of the proof, there exists a function \(g_{i}\in \mathcal{L}^{\infty}(\mathrm{X}_{i},\mu_{i})\) such that \(u_{i} = \widetilde{m}_{g_{i}}\) and \(\mathrm{N}_{\infty}(g_{i})\leqslant \| u_{i}\| \leqslant \| u\|\) .

The function g on X which coincides with \(g_{i}\) on \(X_{i}\) and which is zero on Z is bounded and \(\mu\) -measurable (INT, IV, p. 193, § 5, no. 10, prop. 16), hence \(g \in \mathcal{L}^{\infty}(X, \mu)\) . For every \(i \in I\) , the restriction of u to \(E_{i}\) coincides with that of \(\tilde{m}_{g}\) ; consequently, we have \(u = m_{g}\) by prop. 1, c) of IV, p. 179.

COROLLARY. --- Let u be an endomorphism of the Hilbert space \(\mathrm{L}^{2}(\mathrm{X},\mu)\) commuting with \(\widetilde{m}_{g}\) for every function \(g\in\mathcal{K}(\mathrm{X})\) . Then there exists a unique element \(f\in\mathrm{L}^{\infty}(\mathrm{X},\mu)\) such that \(u=\widetilde{m}_{f}\) .

By prop. 7, it suffices to prove that \(u\) commutes with \(\widetilde{m}_{g}\) for all \(g\in\mathcal{L}^{\infty}(\mathrm{X},\mu)\) . Let \(h_{1}\) and \(h_{2}\) be elements of \(\mathcal{L}^{2}(\mathrm{X},\mu)\) with classes \(\widetilde{h}_{1}\) and \(\widetilde{h}_{2}\) in \(\mathrm{L}^{2}(\mathrm{X},\mu)\) . Let \(k_{1}\) (resp. \(k_{2}\) ) be a function in \(\mathcal{L}^{2}(\mathrm{X},\mu)\) whose class is \(u(\widetilde{h}_{1})\) (resp. \(u^{*}(\widetilde{h}_{2})\) ).

Denote by \(h = h_1 \overline{k}_2 - k_1 \overline{h}_2\) ; we have \(h \in \mathcal{L}^1(\mathrm{X}, \mu)\) . Define the measure \(\nu = h \cdot \mu\) on \(\mathrm{X}\) ; it is bounded. For every \(g \in \mathcal{L}^\infty(\mathrm{X}, \mu)\) , we have

\[
\begin{array}{c} \langle \widetilde {h} _ {2} | u (\widetilde {m} _ {g} (\widetilde {h} _ {1})) - \widetilde {m} _ {g} (u (\widetilde {h} _ {1})) \rangle = \langle u ^ {*} (\widetilde {h} _ {2}) | \widetilde {m} _ {g} (\widetilde {h} _ {1}) \rangle - \langle \widetilde {h} _ {2} | \widetilde {m} _ {g} (u (\widetilde {h} _ {1})) \rangle \\ = \int_ {\mathrm{X}} g \cdot h _ {1} \cdot \overline {{k}} _ {2} d \mu - \int_ {\mathrm{X}} g \cdot k _ {1} \cdot \overline {{h}} _ {2} d \mu = \nu (g). \end{array}
\]

We therefore have by hypothesis \(\nu(g) = 0\) for every function \(g \in \mathcal{K}(\mathrm{X})\) , that is, \(\nu = 0\) . Consequently, it follows that

\[
\langle \widetilde {h} _ {2} | u (\widetilde {m} _ {g} (\widetilde {h} _ {1})) - \widetilde {m} _ {g} (u (\widetilde {h} _ {1})) \rangle = \nu (g) = 0
\]

for all \(g \in \mathcal{L}^{\infty}(\mathrm{X},\mu)\) . Since this is the case for all \(h_1\) and \(h_2\) in \(\mathcal{L}^2 (\mathrm{X},\mu)\) , we have \(u\circ \widetilde{m}_g = \widetilde{m}_g\circ u\) .

Remark. --- In what follows, we will often simply denote by \(m_{g}\) the multiplication operator by g on \(\mathrm{L}^{p}(\mathrm{X},\mu)\) .

